\chapter{Hints and selected solutions}
\label{chap:solutions}

The problems are intended to be worked with the conventions and canonical
definitions of the corresponding Stage.  The following hints give the key
step or a check; they do not replace the intermediate algebra.  A companion
volume, built from \texttt{solutions\_ja.tex}, gives complete Japanese
solutions to all ninety-nine Stage problems and all seven problems in the
resummation appendix.  The present chapter remains intentionally compact so
that it can be used while reading the English text without duplicating that
companion volume.

\section*{Stage I}
For I.2, the boundary form is
$-\rmi\hbar[\phi^*(L)\psi(L)-\phi^*(0)\psi(0)]$; the same phase condition on
$\phi$ and $\psi$ makes it vanish and exhausts the adjoint domain.  For I.4,
$p^2\rmd p=\mu p\rmd E$, so multiplication by $(\mu p)^{1/2}$ preserves the
norm.  For I.8, solve the $Q$ equation with the chosen outgoing boundary value
before substitution; differentiation acts both on the eigenvector and on the
resolvent inside $\vsub{H}{eff}$.
For I.11, approximate the image of the unit ball under the compact operator
by a finite $\varepsilon$-net and use strong convergence on that finite set.
For I.12, square integrability of the kernel is precisely the
Hilbert--Schmidt criterion; inspect the diagonal singularity locally before
estimating the tails.  For I.13, move $V^{1/2}$ through the eigenvalue
equation and keep the sign convention in the Birman--Schwinger operator.
For I.14, insert a Laurent ansatz in $A(z)^{-1}A(z)=\identity$ and pair its
leading equation with the left and right null vectors.  For I.15, choose
normalized functions with mutually disjoint supports: a nonzero multiplier
cannot turn that sequence into a norm-convergent one.

\section*{Stage II}
For II.3, Fourier transform $(E-p^2/2\mu+\rmi0)^{-1}$ and close the radial
momentum contour on the outgoing pole.  For II.6, the radial Green kernel is a
product of the solution regular at the origin and the outgoing solution,
divided by their Wronskian.  For II.7--II.8, the divergent phase is the time
integral of $1/|\bm v t|$ and is logarithmic; the comparison dynamics must
carry exactly that term.
For II.11, integrate the differential expression against a test function
twice and retain the endpoint Wronskian.  For II.12, the two Frobenius powers
are $r^{1/2\pm\sqrt{g+1/4}}$; testing both in $L^2(0,1)$ locates the
limit-point threshold at $g=3/4$.  For II.13, continuity fixes one Green
kernel coefficient and the derivative jump fixes the other.  For II.14,
match the sine solution in the well to the two exterior exponentials and
read the Jost function from the incoming coefficient.  For II.15,
differentiate the differential equation with respect to momentum, subtract
the original equation times the differentiated solution, and integrate the
resulting Wronskian derivative.

\section*{Stage III}
For III.1, equate powers after inserting
$S=\sum_{n=-1}^{\infty}\hbar^nS_n$ into $S^2+S'=Q/\hbar^2$.
For III.2, the Borel transform has a pole at unit Borel time; the two Laplace
contours differ by $2\pi\rmi$ times its residue.  For III.7, a double zero
gives a $(E-E_*)^{-2}$ term unless the numerator vanishes simultaneously.

\section*{Stage IV}
For IV.1, use
$(x\pm\rmi0)^{-1}=\operatorname{PV}(1/x)\mp\rmi\pi\delta(x)$.
For IV.3, write $k=|k|\rme^{-\rmi\varphi}$; the scaled wave decays when
$\theta>\varphi$.  For IV.7, differentiate the Birman--Krein determinant or
take the imaginary part of the reference-subtracted resolvent trace.
For IV.11, expand $D(E)=E-E_0-\Sigma(E)$ at its simple zero; the residue is
$1/D'(\resonantE)=1/[1-\Sigma'(\resonantE)]$.  For IV.12, differentiate the argument of
the partial-wave $S$ matrix before setting $E=\resonantE$.  For IV.13, combine the
background and pole terms over the common denominator
$\epsilon+\rmi$ and complete the modulus square.  For IV.14, put the same
square-root branch into the self-energy, phase space, and analytic
continuation; changing only one of the three changes the sheet.

\section*{Stage V}
For V.3, $u(r)=\rme^{x/2}\chi(x)$ removes the first derivative created by
$r=\rme^x$.  The $1/4$ from this rescaling completes $l(l+1)$ to
$(l+1/2)^2$.  For V.5, evaluate the radial momentum contour by residues at
zero and infinity rather than integrating between turning points directly.

\section*{Stage VI}
For VI.3, step functions approximate every fixed $L^2$ function, but a unit
vector supported inside a newly resolved sub-bin keeps the operator-norm
difference equal to one.  For VI.6, the closed-channel resolvent is real below
threshold and acquires a discontinuity above it.  For VI.10, use
$j_l(qr)\sim(qr)^l/(2l+1)!!$.
For VI.11, write both Jacobi sets explicitly and transform the exchange
operator before imposing the Pauli projector.  For VI.12, the full extended
completeness relation restores the positive response even though individual
rotated-continuum terms need not be positive.  For VI.13, take the outer
product of the right and left pole vectors and divide by their biorthogonal
overlap; a fragmented cluster is treated by a contour integral instead.
For VI.14, begin from the invariant three-body phase-space measure and only
then change variables to the measured energy pair.  For VI.15, if
$\Cincoming(\resonantE)=0$ is simple, solve the connection problem for right and
left null vectors and divide their outer product by
$v_L^\top\Cincoming'(\resonantE)v_R$.

\section*{Stage VII}
For VII.1, the discriminant is
$(\operatorname{tr}H)^2-4\det H$; an EP also requires defectiveness, not only
equal numerical eigenvalues.  For VII.4, one circuit exchanges sheets and two
circuits return the eigenvalue labels while leaving a possible geometric
factor.  For VII.10, solve the determinant and its energy derivative
simultaneously and verify a rank-one null space.

\section*{Stage VIII}
For VIII.1, integrate $\rmd r_*/\rmd r=(1-2M/r)^{-1}$ and choose the additive
constant only after specifying a dimensionless logarithm.  For VIII.5,
$\rme^{-\rmi\omega(t+r_*)}$ is ingoing at the future horizon and
$\rme^{-\rmi\omega(t-r_*)}$ outgoing at infinity.  For VIII.10, retain the
reference-subtracted rotated continuum when reconstructing the response; a
single Lorentzian pole is not the full level density.

\section*{Worked solutions as calculation templates}

The solutions below are deliberately more explicit than the preceding hints.
Each one isolates a calculation that recurs in later Stages.  The point is not
to memorize the final formula, but to make the domain, boundary value, sheet,
or projection choice visible at the line where it enters.

\subsection*{I.2: self-adjoint momentum on an interval}

Begin with the differential expression $-\rmi\hbar\rmd/\rmd x$ on
$L^2(0,L)$ and the domain
\begin{equation}
  \mathcal D(p_\alpha)=
  \{\psi\in H^1(0,L):\psi(L)=\rme^{\rmi\alpha}\psi(0)\}.
\end{equation}
For $\phi,\psi\in H^1(0,L)$, integration by parts gives the boundary form
\begin{align}
  \dualpair{\phi}{p\psi}-\dualpair{p\phi}{\psi}
  &=-\rmi\hbar
  \left[\phi^*(L)\psi(L)-\phi^*(0)\psi(0)\right].
\end{align}
It vanishes on $\mathcal D(p_\alpha)$.  Conversely, if $\phi$ belongs to the
adjoint domain, the boundary form must vanish for every admissible $\psi$.
Using $\psi(L)=\rme^{\rmi\alpha}\psi(0)$ and varying $\psi(0)$ freely yields
$\phi(L)=\rme^{\rmi\alpha}\phi(0)$.  Thus the adjoint has the same domain and
$p_\alpha$ is self-adjoint.

The eigenvalue equation gives $\psi_k(x)=C\rme^{\rmi kx}$.  Its boundary
condition is $\rme^{\rmi kL}=\rme^{\rmi\alpha}$, hence
\begin{equation}
  k_n=\frac{2\pi n+\alpha}{L},
  \qquad
  p_n=\hbar k_n,
  \qquad n\in\mathbb Z .
\end{equation}
The calculation is a compact model of a rule used throughout the book: a
differential expression has no spectrum until its domain has been fixed.

\subsection*{I.8: Feshbach reduction and the derivative metric}

Let $P+Q=\identity$, $PQ=0$, and write $(E-H)\Psi=0$ in blocks:
\begin{align}
  (E-H_{PP})P\Psi-H_{PQ}Q\Psi&=0,
  \\
  -H_{QP}P\Psi+(E-H_{QQ})Q\Psi&=0.
\end{align}
With the specified outgoing boundary value of the $Q$ resolvent,
\begin{equation}
  Q\Psi=(E-H_{QQ})^{-1}H_{QP}P\Psi .
\end{equation}
Substitution gives
\begin{equation}
  [E-\vsub{H}{eff}(E)]P\Psi=0,
  \qquad
  \vsub{H}{eff}(E)=H_{PP}+H_{PQ}(E-H_{QQ})^{-1}H_{QP}.
\end{equation}
For a rank-one $P$ space write the secular function as
$F(E)=E-E_0-\Sigma(E)$.  At a simple pole $\resonantE$,
\begin{equation}
  \frac{1}{F(E)}=
  \frac{Z_R}{E-\resonantE}+O(1),
  \qquad
  Z_R=\frac{1}{1-\Sigma'(\resonantE)}.
\end{equation}
The same factor follows before taking rank one.  Differentiate
$[E-\vsub{H}{eff}(E)]\psi_E=0$, pair with the left pole vector, and evaluate at
$\resonantE$.  The coefficient of the pole is controlled by
$\bra{\widetilde\psi_R}[\identity-\partial_E\vsub{H}{eff}(\resonantE)]\ket{\psi_R}$.
Dropping the derivative would normalize the projected vector as if the
eliminated continuum carried no probability.

\subsection*{II.3: outgoing free Green function and the Born amplitude}

For $E=\hbar^2k^2/(2\mu)$, the outgoing free kernel is
\begin{equation}
  G_0^{(+)}(\bm r;E)
  =\int\frac{\rmd^3p}{(2\pi\hbar)^3}
  \frac{\rme^{\rmi\bm p\cdot\bm r/\hbar}}
  {E-p^2/(2\mu)+\rmi0}.
\end{equation}
After the angular integral, extend the radial momentum integral to the whole
line.  For $r>0$ close the contour in the upper half-plane for the factor
$\rme^{\rmi pr/\hbar}$.  The $+\rmi0$ prescription selects the outgoing pole,
and the residue gives
\begin{equation}
  G_0^{(+)}(\bm r;E)
  =-\frac{\mu}{2\pi\hbar^2}\frac{\rme^{\rmi kr}}{r}.
\end{equation}
Insert this kernel in
$\psi^{(+)}=\rme^{\rmi\bm k\cdot\bm r}+G_0^{(+)}V\psi^{(+)}$.
For $r\gg r'$ use
$|\bm r-\bm r'|=r-\widehat{\bm r}\cdot\bm r'+O(r^{-1})$.
The coefficient of $\rme^{\rmi kr}/r$ is
\begin{equation}
  f(\bm k',\bm k)
  =-\frac{\mu}{2\pi\hbar^2}
  \int\rmd^3r'\,
  \rme^{-\rmi\bm k'\cdot\bm r'}V(\bm r')\psi^{(+)}(\bm r').
\end{equation}
Replacing $\psi^{(+)}$ by the incident plane wave gives the first Born
amplitude.  The sign is therefore fixed by the time convention and the
boundary value of the resolvent, not by a remembered scattering formula.

\subsection*{II.6: Jost Green kernel and pole residue}

For a fixed partial wave let $\varphi_l(k,r)$ be regular at the origin and
$f_l^{(+)}(k,r)$ outgoing at infinity.  Their Wronskian is independent of
$r$.  Solving the jump condition of the radial Green equation gives
\begin{equation}
  G_l^{(+)}(r,r';k)
  =-\frac{2\mu}{\hbar^2}
  \frac{\varphi_l(k,r_<)f_l^{(+)}(k,r_>)}
       {W[\varphi_l,f_l^{(+)}](k)}.
\end{equation}
Up to the convention-dependent nonzero factor, the denominator is the Jost
function $F_l(k)$.  If $F_l(k_R)=0$ and $F_l'(k_R)\ne0$, regular and outgoing
solutions become proportional, say $\varphi_l(k_R,r)=A_Ru_R(r)$ and
$f_l^{(+)}(k_R,r)=B_Ru_R(r)$.  Therefore
\begin{equation}
  \Res_{k=k_R}G_l^{(+)}(r,r';k)
  =-\frac{2\mu}{\hbar^2}
  \frac{A_RB_R}{F_l'(k_R)}u_R(r)u_R(r').
\end{equation}
This is the coordinate-space version of ``pole position plus residue.''  The
state $u_R$ alone is normalization dependent; the full residue is not.

\subsection*{III.1: the exact-WKB Riccati recursion}

Write the Schr"odinger equation as
$[-\hbar^2\partial_x^2+Q(x)]\psi=0$ and set
$\psi=\rme^{\int^xS(x',\hbar)\rmd x'}$.  Then
\begin{equation}
  S^2+\partial_xS=\frac{Q}{\hbar^2}.
\end{equation}
With $S=\sum_{n=-1}^{\infty}\hbar^nS_n$, the first equations are
\begin{align}
  S_{-1}^2&=Q,
  &2S_{-1}S_0+S_{-1}'&=0,
  \\
  2S_{-1}S_1+S_0^2+S_0'&=0,
  &2S_{-1}S_2+2S_0S_1+S_1'&=0.
\end{align}
Choosing $S_{-1}^{(\pm)}=\pm\sqrt Q$ gives
\begin{align}
  S_0&=-\frac{Q'}{4Q},
  \\
  S_1^{(\pm)}&=\pm
  \left(\frac{Q''}{8Q^{3/2}}
  -\frac{5(Q')^2}{32Q^{5/2}}\right),
  \\
  S_2&=-\frac{Q'''}{16Q^2}
  +\frac{9Q'Q''}{32Q^3}
  -\frac{15(Q')^3}{64Q^4}.
\end{align}
The coefficients that change sign with the square-root sheet form
$\vsub{S}{odd}$; the sheet-even coefficients form $\vsub{S}{even}$, with
$\vsub{S}{even}=-(1/2)\partial_x\log \vsub{S}{odd}$.  This identity fixes the
all-order prefactor and provides a useful algebraic check on any recursion
code.

\subsection*{III.7: simple and double zeros of a connection coefficient}

Let $S(E)=N(E)/\Cincoming(E)$ and put $z=E-E_*$.  At a simple zero,
\begin{align}
  \Cincoming(E)&=c_1z+\tfrac12c_2z^2+O(z^3),
  &N(E)&=n_0+n_1z+O(z^2).
\end{align}
Long division gives
\begin{equation}
  S(E)=\frac{n_0/c_1}{z}
  +\left(\frac{n_1}{c_1}-\frac{n_0c_2}{2c_1^2}\right)+O(z).
\end{equation}
At a double zero, $c_1=0$ and $c_2\ne0$.  Including the cubic coefficient,
\begin{equation}
  S(E)=\frac{2n_0/c_2}{z^2}
  +\frac{2}{c_2}
   \left(n_1-\frac{n_0c_3}{3c_2}\right)\frac{1}{z}+O(1).
\end{equation}
The second-order pole is the scalar shadow of a length-two Jordan chain.
Notice the necessary caveat: if $N$ vanishes at the same point, cancellation
can lower the pole order.  A multiple zero of a numerical determinant must
therefore be checked against the relevant numerator and nullity.

\subsection*{IV.1: a flat-continuum self-energy}

Take one bare state of energy $E_0$, continuum energies
$\epsilon\in[\vsub{E}{th},E_c]$, constant density $\rho$, and constant coupling
$g$.  The outgoing self-energy is
\begin{equation}
  \Sigma(E+\rmi0)=\rho|g|^2
  \int_{\vsub{E}{th}}^{E_c}
  \frac{\rmd\epsilon}{E-\epsilon+\rmi0}.
\end{equation}
Using the distribution identity before performing the integral,
\begin{align}
  \Sigma(E+\rmi0)
  &=\rho|g|^2
  \log\left|\frac{E-\vsub{E}{th}}{E-E_c}\right|
  -\rmi\pi\rho|g|^2
  \Theta(E-\vsub{E}{th})\Theta(E_c-E).
\end{align}
Thus the real part shifts the level and depends on the ultraviolet model,
whereas inside the band
$\Gamma(E)=-2\im\Sigma(E+\rmi0)=2\pi\rho|g|^2$.
The pole solves $\resonantE-E_0-\vsup{\Sigma}{II}(\resonantE)=0$ on the continued sheet.
Keeping the logarithm rather than replacing it immediately by a constant
width also retains the threshold branch point and the nonexponential tail.

\subsection*{IV.5: Berggren--Zel'dovich Gaussian continuation}

The asymptotic contribution to the Gamow pairing contains
\begin{equation}
  I_\epsilon(q)=\int_0^\infty
  \rme^{-\epsilon r^2+\rmi q r}\rmd r,
  \qquad \re\epsilon>0.
\end{equation}
Complete the square and use the complementary error function:
\begin{equation}
  I_\epsilon(q)=\frac{\sqrt\pi}{2\sqrt\epsilon}
  \rme^{-\frac{q^2}{4\epsilon}}
  \operatorname{erfc}\left(-\frac{\rmi q}{2\sqrt\epsilon}\right).
\end{equation}
In a sector where $\im q>0$, the undamped integral already converges and
$I_\epsilon(q)\to-1/(\rmi q)$ as $\epsilon\downarrow0$.  The right-hand side
is analytic in $q$ away from the Stokes boundary, so this value defines its
continuation to the resonance quadrant even though the original undamped
real-axis integral diverges there.  For a full wave function one performs
this continuation term by term only after controlling the subleading
asymptotics.  Berggren's proof is precisely what prevents ``insert a Gaussian
and discard it'' from becoming a regulator-dependent prescription.

\subsection*{V.3: deriving, not guessing, the Langer term}

For the reduced three-dimensional radial function,
\begin{equation}
  -\hbar^2u''(r)+
  \left[2\mu V(r)+\frac{\hbar^2l(l+1)}{r^2}\right]u(r)
  =2\mu Eu(r).
\end{equation}
Set $r=\rme^x$ and $u(r)=\rme^{x/2}\chi(x)$.  Since
$\partial_r=\rme^{-x}\partial_x$,
\begin{equation}
  u''(r)=\rme^{-3x/2}
  \left(\chi''-\frac14\chi\right).
\end{equation}
Multiplying the equation by $\rme^{3x/2}$ produces
\begin{equation}
  -\hbar^2\chi''(x)+
  \left\{2\mu\rme^{2x}[V(\rme^x)-E]
  +\hbar^2\left(l+\frac12\right)^2\right\}\chi(x)=0.
\end{equation}
The extra $1/4$ is the Jacobian contribution required to remove the first
derivative.  The replacement $l(l+1)\mapsto(l+1/2)^2$ is therefore a
coordinate-and-half-density statement; using it without transforming the
wave function mixes two different differential equations.

\subsection*{V.5: the radial Coulomb action}

For $V(r)=-\kappa/r$ and $E<0$, put
$K=\sqrt{-2\mu E}$ and $L=\hbar(l+1/2)$.  The Langer radial momentum is
\begin{equation}
  p_r(r)=\sqrt{-K^2+\frac{2\mu\kappa}{r}-\frac{L^2}{r^2}}.
\end{equation}
The closed contour around the two positive turning points can be deformed to
small and large circles on the two-sheeted curve.  Expanding at zero and
infinity, or evaluating the corresponding residues, yields
\begin{equation}
  \frac{1}{2\pi}\oint p_r\rmd r
  =\frac{\mu\kappa}{K}-L.
\end{equation}
Bohr--Sommerfeld quantization of the radial oscillation is
$(2\pi)^{-1}\oint p_r\rmd r=\hbar(n_r+1/2)$.  Hence
\begin{equation}
  \frac{\mu\kappa}{K}=\hbar(n_r+l+1),
  \qquad
  E=-\frac{\mu\kappa^2}
  {2\hbar^2(n_r+l+1)^2}.
\end{equation}
Keeping the origin residue $-L$ and the turning-point half phase separate
shows how the Frobenius monodromy and ordinary turning points cooperate.

\subsection*{VI.2: projection to the CDCC equations}

Let $\bm r$ be the projectile internal coordinate and $\bm R$ its
target separation.  Expand the total wave function in retained projectile
states and channel angular functions,
\begin{equation}
  \Psi^{JM}(\bm R,\bm r)
  =\sum_c\frac{\chi_c^J(R)}{R}
  \left[\phi_{nI}(\bm r)\otimes Y_L(\widehat{\bm R})\right]_{JM},
  \qquad c=(n,I,L).
\end{equation}
Insert this expansion in
$[T_R+h(\bm r)+U_{aA}+U_{bA}-E]\Psi=0$, multiply by the conjugate channel
function, and integrate over $\bm r$ and $\widehat{\bm R}$.  Orthogonality and
$h\phi_{nI}=\epsilon_{nI}\phi_{nI}$ give
\begin{equation}
  \left[-\frac{\hbar^2}{2\mu_R}
  \left(\frac{\rmd^2}{\rmd R^2}
  -\frac{L_c(L_c+1)}{R^2}\right)
  +\epsilon_c-E\right]\chi_c^J(R)
  +\sum_{c'}U^J_{cc'}(R)\chi_{c'}^J(R)=0,
\end{equation}
where
\begin{equation}
  U^J_{cc'}(R)=
  \Bra{[\phi_{nI}\otimes Y_L]_{JM}}
  U_{aA}(\bm R_a)+U_{bA}(\bm R_b)
  \Ket{[\phi_{n'I'}\otimes Y_{L'}]_{JM}} .
\end{equation}
Nothing in this derivation makes a positive-energy pseudostate a bound
state.  It is a basis vector in the finite model-space projection; physical
breakup information is recovered only after convergence and smoothing.

\subsection*{VI.6: eliminating a closed reaction channel}

For two coupled radial channels write
\begin{align}
  (E-H_{11})\chi_1&=V_{12}\chi_2,
  &(E-H_{22})\chi_2&=V_{21}\chi_1.
\end{align}
The boundary condition in channel 2 fixes its Green operator, so
\begin{equation}
  \left[E-H_{11}-\vsub{U}{pol}(E)\right]\chi_1=0,
  \qquad
  \vsub{U}{pol}(E)=V_{12}(E-H_{22})^{-1}V_{21}.
\end{equation}
Below the second threshold, a self-adjoint $H_{22}$ has no on-shell spectral
weight at $E$.  Its resolvent boundary value is real, and $\vsub{U}{pol}$ shifts
and distorts the elastic potential without absorbing flux.  Above threshold,
Stone's formula gives
\begin{equation}
  \im \vsub{U}{pol}(E+\rmi0)
  =-\pi V_{12}\delta(E-H_{22})V_{21},
\end{equation}
which is negative semidefinite and supplies the reaction loss.  Omitting
closed CDCC channels can nevertheless change an elastic observable below
threshold because their real dispersive contribution remains.

\subsection*{VII.1: the algebraic test for a second-order EP}

For
\begin{equation}
  H=\begin{pmatrix}a&b\\b&d\end{pmatrix},
  \qquad a,b,d\in\mathbb C,
\end{equation}
the eigenvalues are
\begin{equation}
  \lambda_\pm=\frac{a+d}{2}
  \pm\frac12\sqrt{(a-d)^2+4b^2}.
\end{equation}
They coalesce when the discriminant
$D=(a-d)^2+4b^2$ vanishes.  This condition is necessary but must be joined to
a rank test.  If $H=\lambda_*\identity$, the eigenspace is two-dimensional and
the degeneracy is ordinary.  Otherwise
$\rank(H-\lambda_*\identity)=1$, the nullity is one, and the algebraic
multiplicity two exceeds the geometric multiplicity: the point is an EP.
Near a generic control value $g_*$, $D(g)=D'(g_*)(g-g_*)+\cdots$, so
\begin{equation}
  \lambda_+-\lambda_-
  \sim\sqrt{D'(g_*)}\,(g-g_*)^{1/2}.
\end{equation}
The square-root exchange is thus derived from the secular equation.  A plot
showing close eigenvalues is not an EP diagnostic until defectiveness and the
local Puiseux law have also been checked.

\subsection*{VII.2: resolvent of a length-two Jordan chain}

Choose right vectors satisfying
\begin{equation}
  (H-E_*)\psi_0=0,
  \qquad
  (H-E_*)\psi_1=\psi_0,
\end{equation}
and a dual chain normalized so that the nilpotent part is represented by
$N=\ket{\psi_0}\bra{\widetilde\psi_1}$ with $N^2=0$ on the root
subspace.  There $H=E_*\identity+N$.  For $z\ne E_*$,
\begin{align}
  (z-H)^{-1}
  &=\left[(z-E_*)\identity-N\right]^{-1}
  \\
  &=\frac{\identity}{z-E_*}
  +\frac{N}{(z-E_*)^2},
\end{align}
because the geometric series terminates at $N^2=0$.  Restoring the explicit
root-space identity gives a simple-pole projector assembled from both chain
vectors and a double-pole coefficient proportional to
$\ket{\psi_0}\bra{\widetilde\psi_1}$.  This is why solving only for the
coalesced eigenvector loses information: the associated vector controls the
principal part of the response.

\subsection*{VIII.1: Schwarzschild tortoise coordinate}

For $f(r)=1-2M/r$, radial null curves suggest defining
$\rmd r_*/\rmd r=f^{-1}$.  Partial fractions give
\begin{equation}
  r_*=r+2M\log\left|\frac{r}{2M}-1\right|+C.
\end{equation}
Outside the horizon, $r>2M$.  As $r=2M+\varepsilon$ with
$\varepsilon\downarrow0$,
\begin{equation}
  r_*=2M\log\frac{\varepsilon}{2M}+O(1)\longrightarrow-\infty.
\end{equation}
At large radius,
$r_*=r+2M\log(r/2M)+O(1)\to+\infty$.  The logarithm's argument has been made
dimensionless; changing its reference scale merely changes $C$.  Under
complex continuation, however, the logarithm also records a branch choice.
That branch structure is why an unrestricted rotation of the entire
$r_*$ plane cannot be assumed to induce a single-valued rotation of $r$.

\subsection*{VIII.5: horizon signs as connection data}

With time dependence $\rme^{-\rmi\omega t}$, introduce
$v=t+r_*$ and $u=t-r_*$.  Near the future event horizon a smooth ingoing wave
depends on $v$:
\begin{equation}
  \rme^{-\rmi\omega v}
  =\rme^{-\rmi\omega t}\rme^{-\rmi\omega r_*}.
\end{equation}
At spatial infinity an outgoing wave depends on $u$:
\begin{equation}
  \rme^{-\rmi\omega u}
  =\rme^{-\rmi\omega t}\rme^{+\rmi\omega r_*}.
\end{equation}
Let $\vsupb{f}{in}{H}$ be the horizon-normalized ingoing solution.  At
infinity it has the connection expansion
\begin{equation}
  \vsupb{f}{in}{H}
  =\vsub{\mathcal C}{out}(\omega)\vsup{f_\infty}{out}
  +\Cincoming(\omega)\vsup{f_\infty}{in}.
\end{equation}
A QNM satisfies both physical conditions exactly when
$\Cincoming(\omega)=0$.  For $\im\omega<0$ these exponentials diverge
on parts of the real $r_*$ axis; that divergence is not a sign error but the
reason analytic continuation, complex scaling, or a rigged-space pairing is
needed to represent the same connection zero.
